\ifdefined\gkver\else\def\gkver{student}\fi
\documentclass[\gkver]{gaokaozhenti}
\begin{document}

\chapter{2016年浙江理综}
%% paperType: 理综物理部分

\begin{enumerate}
\item
%% number: 14
%% paperName: 2016年浙江理综第 14 题 6 分
%% typeId: 1
%% type: 单选题
%% chapter: 电场
%% point: 电势能电势差电场力做功
%% method: 无
%% score: 6
%% degree: 750
%% duplicateId: 0
%% body:
以下说法正确的是\xzanswer{A}
\fourchoices[answer=A]
{在静电场中，沿着电场线方向电势逐渐降低}
{外力对物体所做的功越多，对应的功率越大}
{电容器电容$C$与电容器所带电荷量$Q$成正比}
{在超重和失重现象中，地球对物体的实际作用力发生了变化}
%% answer: A
%% memo:
\memoanswer{
	本题原卷及材料中未提供详解，待补充。
}

\item
%% number: 15
%% paperName: 2016年浙江理综第 15 题 6 分
%% typeId: 1
%% type: 单选题
%% chapter: 电场
%% point: 静电感应
%% method: 无
%% score: 6
%% degree: 650
%% duplicateId: 0
%% body:
如图所示，两个不带电的导体A和B，用一对绝缘柱支持使它们彼此接触。把一带正电荷的物体C置于A附近，贴在A、B下部的金属箔都张开，\xzanswer{C}
\onepicture[w=3.2cm]{figs/15.png}
\fourchoices[answer=C]
{此时A带正电，B带负电}
{此时A电势低，B电势高}
{移去C，贴在A、B下部的金属箔都闭合}
{先把A和B分开，然后移去C，贴在A、B下部的金属箔都闭合}
%% answer: C
%% memo:
\memoanswer{
	本题原卷及材料中未提供详解，待补充。
}

\item
%% number: 16
%% paperName: 2016年浙江理综第 16 题 6 分
%% typeId: 1
%% type: 单选题
%% chapter: 电磁感应
%% point: 楞次定律
%% method: 无
%% score: 6
%% degree: 550
%% duplicateId: 0
%% body:
如图所示，a、b两个闭合正方形线圈用同样的导线制成，匝数均为10匝，边长$l_{a}=3l_{b}$，图示区域内有垂直纸面向里的匀强磁场，且磁感应强度随时间均匀增大，不考虑线圈之间的相互影响，则\xzanswer{B}
\onepicture[w=3.2cm]{figs/16.png}
\fourchoices[answer=B]
{两线圈内产生顺时针方向的感应电流}
{a、b线圈中感应电动势之比为9:1}
{a、b线圈中感应电流之比为3:4}
{a、b线圈中电功率之比为3:1}
%% answer: B
%% memo:
\memoanswer{
	根据楞次定律可知，两线圈内均产生逆时针方向的感应电流，选项 A 错误；因磁感应强度随时间均匀增大，则 $\frac{\Delta B}{\Delta t}=k$，根据法拉第电磁感应定律可知 $E=n\frac{\Delta \Phi}{\Delta t}=n\frac{\Delta B}{\Delta t}l^{2}$，则 $\frac{E_{a}}{E_{b}}=\left(\frac{3}{1}\right)^{2}=\frac{9}{1}$，选项 B 正确；根据 $I=\frac{E}{R}=\frac{E}{\rho\frac{4nl}{S'}}=\frac{n\frac{\Delta B}{\Delta t}l^{2}S'}{4\rho nl}=\frac{klS'}{4\rho}\propto l$，故 a、b 线圈中感应电流之比为 3:1，选项 C 错误；电功率 $P=IE=\frac{klS'}{4\rho}\cdot n\frac{\Delta B}{\Delta t}l^{2}=\frac{nk^{2}l^{3}S'}{4\rho}\propto l^{3}$，故 a、b 线圈中电功率之比为 27:1，选项 D 错误；故选（B）。
}

\item
%% number: 17
%% paperName: 2016年浙江理综第 17 题 6 分
%% typeId: 1
%% type: 单选题
%% chapter: 电路
%% point: 传感器
%% method: 无
%% score: 6
%% degree: 450
%% duplicateId: 0
%% body:
如图所示为一种常见的身高体重测量仪。测量仪顶部向下发射波速为 $v$ 的超声波，超声波经反射后返回，被测量仪接收，测量仪记录发射和接收的时间间隔。质量为 $M_{0}$ 的测重台置于压力传感器上，传感器输出电压与作用在其上的压力成正比。当测重台没有站人时，测量仪记录的时间间隔为 $t_{0}$，输出电压为 $U_{0}$，某同学站上测重台，测量仪记录的时间间隔为 $t$，输出电压为 $U$，则该同学身高和质量分别为\xzanswer{D}
\onepicture[h=3.4cm]{figs/17.png}
\fourchoices[answer=D]
{$v(t_{0}-t)$，$\frac{U}{U_{0}}M_{0}$}
{$\frac{1}{2}v(t_{0}-t)$，$\frac{U}{U_{0}}M_{0}$}
{$v(t_{0}-t)$，$\frac{U-U_{0}}{U_{0}}M_{0}$}
{$\frac{1}{2}v(t_{0}-t)$，$\frac{U-U_{0}}{U_{0}}M_{0}$}
%% answer: D
%% memo:
\memoanswer{
	本题原卷及材料中未提供详解，待补充。
}

\item
%% number: 18
%% paperName: 2016年浙江理综第 18 题 6 分
%% typeId: 2
%% type: 多选题
%% chapter: 功和能
%% point: 动能定理
%% method: 无
%% score: 6
%% degree: 450
%% duplicateId: 0
%% body:
如图所示为一滑草场。某条滑道由上下两段高均为$h$，与水平面倾角分别为$45^\circ$和$37^\circ$的滑道组成，滑草车与草地之间的动摩擦因数为$\mu$。质量为$m$的载人滑草车从坡顶由静止开始自由下滑，经过上、下两段滑道后，最后恰好静止于滑道的底端（不计滑草车在两段滑道交接处的能量损失，$\sin37^\circ=0.6$，$\cos37^\circ=0.8$）。则\xzanswer{AB}
\onepicture[w=4.0cm]{figs/18.png}
\fourchoices[answer=AB]
{动摩擦因数$\mu=\frac{6}{7}$}
{载人滑草车最大速度为$\sqrt{\frac{2gh}{7}}$}
{载人滑草车克服摩擦力做功为$mgh$}
{载人滑草车在下段滑道上的加速度大小为$\frac{3}{5}g$}
%% answer: AB
%% memo:
\memoanswer{
	本题原卷及材料中未提供详解，待补充。
}

\item
%% number: 19
%% paperName: 2016年浙江理综第 19 题 6 分
%% typeId: 2
%% type: 多选题
%% chapter: 电场
%% point: 带电粒子的平衡
%% method: 无
%% score: 6
%% degree: 450
%% duplicateId: 0
%% body:
如图所示，把A、B两个相同的导电小球分别用长为$0.10\Um$的绝缘细线悬挂于$O_{A}$和$O_{B}$两点。用丝绸摩擦过的玻璃棒与A球接触，棒移开后将悬点$O_{B}$移到$O_{A}$点固定。两球接触后分开，平衡时距离为$0.12\Um$。已测得每个小球质量是$8.0\times10^{-4}\Ukg$，带电小球可视为点电荷，重力加速度$g=10\Umsq$，静电力常量$k=9.0\times10^{9}\UNmqCq$，则\xzanswer{ACD}
\onepicture[h=3.6cm]{figs/19.png}
\fourchoices[answer=ACD]
{两球所带电荷量相等}
{A球所受的静电力为$1.0\times10^{-2}\UN$}
{B球所带的电荷量为$4\sqrt{6}\times10^{-8}\UC$}
{A、B两球连线中点处的电场强度为0}
%% answer: ACD
%% memo:
\memoanswer{
	本题原卷及材料中未提供详解，待补充。
}

\item
%% number: 20
%% paperName: 2016年浙江理综第 20 题 6 分
%% typeId: 2
%% type: 多选题
%% chapter: 曲线运动
%% point: 圆周运动的应用
%% method: 无
%% score: 6
%% degree: 450
%% duplicateId: 0
%% body:
如图所示为赛车场的一个水平“梨形”赛道，两个弯道分别为半径$R=90\Um$的大圆弧和$r=40\Um$的小圆弧，直道与弯道相切。大、小圆弧圆心O、O$'$距离$L=100\Um$。赛车沿弯道路线行驶时，路面对轮胎的最大径向静摩擦力是赛车重力的2.25倍。假设赛车在直道上做匀变速直线运动，在弯道上做匀速圆周运动，要使赛车不打滑，绕赛道一圈时间最短（发动机功率足够大，重力加速度$g=10\Umsq$，$\pi=3.14$），则赛车\xzanswer{AB}
\onepicture[w=3.6cm]{figs/20.png}
\fourchoices[answer=AB]
{在绕过小圆弧弯道后加速}
{在大圆弧弯道上的速率为$45\Ums$}
{在直道上的加速度大小为$5.63\Umsq$}
{通过小圆弧弯道的时间为$5.85\Us$}
%% answer: AB
%% memo:
\memoanswer{
	本题原卷及材料中未提供详解，待补充。
}

\item
%% number: 21
%% paperName: 2016年浙江理综第 21 题 10 分
%% typeId: 4
%% type: 实验题
%% chapter: 力物体的平衡
%% point: 验证平行四边形实验
%% method: 无
%% score: 10
%% degree: 550
%% duplicateId: 0
%% body:
某同学在“探究弹力和弹簧伸长的关系”的实验中，测得图中弹簧OC的劲度系数为$500\UNm$。如图 \ref{2016浙江理综21a} 所示，用弹簧OC和弹簧秤a、b做“探究求合力的方法”实验。在保持弹簧伸长$1.00\Ucm$不变的条件下，
\begin{enumerate}
	\item
	若弹簧秤a、b间夹角为$90^\circ$，弹簧秤a的读数是\tkanswer{3.00$\sim$3.02}$\UN$（图 \ref{2016浙江理综21b} 中所示），则弹簧秤b的读数可能为\tkanswer{3.9$\sim$4.1}$\UN$。

	\item
	若弹簧秤a、b间夹角大于$90^\circ$，保持弹簧秤a与弹簧OC的夹角不变，减小弹簧秤b与弹簧OC的夹角，则弹簧秤a的读数\tkanswer{变大}、弹簧秤b的读数\tkanswer{变大}（填“变大”、“变小”或“不变”）。

\end{enumerate}
\twopicture[align=right, wA=4.0cm, wB=1.8cm, labelA=2016浙江理综21a, labelB=2016浙江理综21b]{figs/21a.png}{figs/21b.png}
%% answer: 
\jdanswer{
\begin{enumerate}
	\item
	3.00$\sim$3.02，3.9$\sim$4.1（有效数不作要求）

	\item
	变大，变大

\end{enumerate}
}
%% memo:
\memoanswer{
	本题原卷及材料中未提供详解，待补充。
}

\item
%% number: 22
%% paperName: 2016年浙江理综第 22 题 10 分
%% typeId: 4
%% type: 实验题
%% chapter: 电路
%% point: 伏安法测电阻
%% method: 无
%% score: 10
%% degree: 450
%% duplicateId: 0
%% body:
某同学用伏安法测量导体的电阻，现有量程为$3\UV$、内阻约为$3\UkO$的电压表和量程为$0.6\UA$、内阻约为$0.1\UO$的电流表。采用分压电路接线，图 \ref{2016浙江理综22a} 是实物的部分连线图，待测电阻为图 \ref{2016浙江理综22b} 中的$R_{1}$，其阻值约为$5\UO$。
\begin{enumerate}
	\item
	测$R_{1}$阻值的最优连接方式为导线①连接\tkanswer{b}（填a或b）、导线②连接\tkanswer{c}（填c或d）。

	\item
	正确接线测得实验数据如表，用作图法求得$R_{1}$的阻值为\tkanswer{4.4}$\UO$。

	\begin{center}
	\begin{tabular}{|c|c|c|c|c|c|c|}
	\hline
	$U$/A & 0.40 & 0.80 & 1.20 & 1.60 & 2.00 & 2.40 \\
	\hline
	$I$/A & 0.09 & 0.19 & 0.27 & 0.35 & 0.44 & 0.53 \\
	\hline
	\end{tabular}
	\end{center}

	\item
	已知图 \ref{2016浙江理综22b} 中$R_{2}$与$R_{1}$是材料相同、厚度相等、表面为正方形的两导体，$R_{2}$的边长是$R_{1}$的$\frac{1}{10}$，若测$R_{2}$的阻值，则最优的连线应选\tkanswer{C}（填选项）。

	（A）①连接a，②连接c \qquad （B）①连接a，②连接d

	（C）①连接b，②连接c \qquad （D）①连接b，②连接d

\end{enumerate}
\twopicture[align=right, wA=3.4cm, wB=2.0cm, labelA=2016浙江理综22a, labelB=2016浙江理综22b]{figs/22a.png}{figs/22b.png}
%% answer: 
\jdanswer{
\begin{enumerate}
	\item
	b，c

	\item
	4.4（4.4$\sim$4.5均可）

	\item
	C

\end{enumerate}
}
%% memo:
\memoanswer{
	本题原卷及材料中未提供详解，待补充。
}

\item
%% number: 23
%% paperName: 2016年浙江理综第 23 题 16 分
%% typeId: 6
%% type: 计算题
%% chapter: 曲线运动
%% point: 平抛运动
%% method: 无
%% score: 16
%% degree: 550
%% duplicateId: 0
%% body:
在真空环境内探测微粒在重力场中能量的简化装置如图所示。P是一个微粒源，能持续水平向右发射质量相同、初速度不同的微粒。高度为$h$的探测屏AB竖直放置，离P点的水平距离为$L$，上端A与P点的高度差也为$h$。
\begin{enumerate}
	\item
	若微粒打在探测屏AB的中点，求微粒在空中飞行的时间；

	\item
	求能被屏探测到的微粒的初速度范围；

	\item
	若打在探测屏A、B两点的微粒的动能相等，求$L$与$h$的关系。

\end{enumerate}
\onepicture[w=3.0cm, align=right]{figs/23.png}
%% answer: 
\jdanswer{
\begin{enumerate}
	\item
	$t=\sqrt{\frac{3h}{g}}$

	\item
	$L\sqrt{\frac{g}{4h}}\leq v\leq L\sqrt{\frac{g}{2h}}$

	\item
	$L=2\sqrt{2}h$

\end{enumerate}
}
%% memo:
\memoanswer{
	本题原卷及材料中未提供详解，待补充。
}

\item
%% number: 24
%% paperName: 2016年浙江理综第 24 题 20 分
%% typeId: 6
%% type: 计算题
%% chapter: 电磁感应
%% point: 电磁感应受力分析
%% method: 无
%% score: 20
%% degree: 550
%% duplicateId: 0
%% body:
小明设计的电磁健身器的简化装置如图所示，两根平行金属导轨相距$l=0.50\Um$，倾角$\theta=53^\circ$，导轨上端串接一个$R=0.05\UO$的电阻。在导轨间长$d=0.56\Um$的区域内，存在方向垂直导轨平面向下的匀强磁场，磁感应强度$B=2.0\UT$。质量$m=4.0\Ukg$的金属棒CD水平置于导轨上，用绝缘绳索通过定滑轮与拉杆GH相连。CD棒的初始位置与磁场区域的下边界相距$s=0.24\Um$。一位健身者用恒力$F=80\UN$拉动GH杆，CD棒由静止开始运动，上升过程中CD棒始终保持与导轨垂直。当CD棒到达磁场上边界时健身者松手，触发恢复装置使CD棒回到初始位置（重力加速度$g=10\Umsq$，$\sin 53^\circ=0.8$，不计其他电阻、摩擦力以及拉杆和绳索的质量）。求：
\begin{enumerate}
	\item
	CD棒进入磁场时速度$v$的大小；

	\item
	CD棒进入磁场时所受的安培力$F_{A}$的大小；

	\item
	在拉升CD棒的过程中，健身者所做的功$W$和电阻产生的焦耳热$Q$。

\end{enumerate}
\onepicture[w=4.0cm, align=right]{figs/24.png}
%% answer: 
\jdanswer{
\begin{enumerate}
	\item
	$2.4\Ums$

	\item
	$48\UN$

	\item
	$64\UJ$，$26.88\UJ$

\end{enumerate}
}
%% memo:
\memoanswer{
	本题原卷及材料中未提供详解，待补充。
}

\item
%% number: 25
%% paperName: 2016年浙江理综第 25 题 22 分
%% typeId: 6
%% type: 计算题
%% chapter: 磁场
%% point: 带电粒子在磁场中的运动
%% method: 无
%% score: 22
%% degree: 350
%% duplicateId: 0
%% body:
为了进一步提高回旋加速器的能量，科学家建造了“扇形聚焦回旋加速器”。在扇形聚焦过程中，离子能以不变的速率在闭合平衡轨道上周期性旋转。

扇形聚焦磁场分布的简化图如图所示，圆心为O的圆形区域等分成六个扇形区域，其中三个为峰区，三个为谷区，峰区和谷区相间分布。峰区内存在方向垂直纸面向里的匀强磁场，磁感应强度为$B$，谷区内没有磁场。质量为$m$，电荷量为$q$的正离子，以不变的速率$v$旋转，其闭合平衡轨道如图中虚线所示。
\begin{enumerate}
	\item
	求闭合平衡轨道在峰区内圆弧的半径$r$，并判断离子旋转的方向是顺时针还是逆时针；

	\item
	求轨道在一个峰区内圆弧的圆心角$\theta$，及离子绕闭合平衡轨道旋转的周期$T$；

	\item
	在谷区也施加垂直纸面向里的匀强磁场，磁感应强度为$B'$，新的闭合平衡轨道在一个峰区内的圆心角$\theta$变为$90^\circ$，求$B'$和$B$的关系。已知：$\sin(\alpha\pm\beta)=\sin\alpha\cos\beta\pm\cos\alpha\sin\beta$，$\cos\alpha=1-2\sin^{2}\frac{\alpha}{2}$

\end{enumerate}
\onepicture[w=3.4cm, align=right]{figs/25.png}
%% answer: 
\jdanswer{
\begin{enumerate}
	\item
	$\frac{mv}{qB}$，旋转方向为逆时针方向

	\item
	$\theta=\frac{2\pi}{3}$，$T=\frac{(2\pi+3\sqrt{3})m}{qB}$

	\item
	$B'=\frac{\sqrt{3}-1}{2}B$

\end{enumerate}
}
%% memo:
\memoanswer{
	本题原卷及材料中未提供详解，待补充。
}

\end{enumerate}

\end{document}
